\documentclass[10pt,a4paper]{amsart}
\usepackage{amsmath,amssymb,amsthm,graphicx}
\usepackage{hyperref}
% Minimal class-compatibility preamble for a reader excerpt of
%   Maxim Nazarov, "Yangian of the General Linear Lie Superalgebra",
%   SIGMA 16 (2020), 112, DOI 10.3842/SIGMA.2020.112.
%
% The block below is the author's own macro and theorem declaration block of the pinned
% source file, copied mechanically and unchanged (source lines 9-71). The SIGMA publisher
% class sigma.cls is not installed here, so the wrapper class declares the same theorem
% environments; no mathematical macro, statement, symbol or assumption is introduced,
% changed or removed by this file.
\numberwithin{equation}{section}

\newtheorem{Theorem}{Theorem}[section]
\newtheorem{Corollary}[Theorem]{Corollary}
\newtheorem{Lemma}[Theorem]{Lemma}
\newtheorem{Proposition}[Theorem]{Proposition}

\def\1{^{-1}}

\def\a{\mathfrak{a}}
\def\al{\alpha}

\def\be{\beta}
\def\bi{{\bar\imath }}
\def\bj{{\bar\jmath }}
\def\bk{{\bar k }}
\def\bl{{\bar l }}
\def\bp{{\bar p }}
\def\bq{{\bar q }}
\def\br{{\bar r }}
\def\bs{{\bar s }}

\def\CC{{\mathbb{C}}}
\def\CMN{\CC^{M|N}}

\def\de{\delta}
\def\De{\Delta}

\def\End{\operatorname{End}}
\def\ep{\varepsilon}

\def\g{\glMN [u]}
\def\ge{\geqslant}
\def\glMN{\mathfrak{gl}_{M|N}}
\def\gr{{\rm gr} }
\def\grpr{{\rm gr}^{\prime} }

\def\id{{\mathrm{id}}}

\def\le{\leqslant}

\def\ns{{\hskip-.5pt}}

\def\om{\omega}
\def\ot{\otimes}

\def\RS{R^{\sharp}}

\def\S{\operatorname{S}}
\def\si{\sigma}
\def\str{\operatorname{str}}
\def\Sym{\mathfrak{S}}

\def\TD{\dot{T}}
\def\TP{T^{\prime}}
\def\TS{T^{\sharp}}

\def\Ua{\operatorname{U}(\a [u])}
\def\UMN{\operatorname{U}(\glMN)}
\def\Ug{\operatorname{U}(\g)}
\def\XP{X^{\prime}}
\def\YMN{\operatorname{Y}(\glMN)}
\def\ZZ{\mathbb{Z}}
\title{Existence of the central series Z(u) of the general linear super Yangian with its two RTT-type identities}
\author{Maxim Nazarov}
\date{doi:10.3842/SIGMA.2020.112; CC BY 4.0}
\begin{document}
\maketitle
\setcounter{section}{1}
\noindent\textit{Editorial note.} The counted unit is Proposition 3.1 at source lines 877-886, source label zu, with the proof environment that begins on the next line at 888-977; the pairing is the ordinary adjacent one. The printed number is fixed by initial\_counters so that every retained item keeps its original number: the source has the Introduction as section 1, so starting the counter at 1 makes the retained Definition of the Yangian section 2 (with its Proposition 2.1) and the Centre of the Yangian section 3, where the counted proposition prints as Proposition 3.1 exactly as in the published article.\par\medskip
\noindent\textit{Editorial note.} The retained definition passages at source lines 169-351 and 368-422 are the author's own Section 1 on the definition of the Yangian: the graded tensor conventions, the generators and their defining relations, the permutation operator P, the R-matrix with (3.52), the series T(u) and T\_p(u), Propositions 2.1 and 2.2 with their proofs, the antiautomorphism tau of End C\textasciicircum{}\{M|N\} defined at line 368, Proposition 2.3 with its proof, and the series T\textasciicircum{}S(u) at line 414. The counted proof uses tau in its first line through the operator Q, so both passages are retained; they carry the section heading of the source and are context, not a second counted proposition.\par\medskip
\noindent\textit{Editorial note.} The context span at 863-876 is the opening of the Centre of the Yangian section, including the author's remark that the description of the centre comes from computing the square of the antipodal mapping. No retained passage was altered.\par\medskip
\noindent\textit{Editorial note.} The retained preamble reproduces the author's macro and theorem declaration block verbatim from source lines 9-71, including the same theorem environment names the SIGMA publisher class declares; the publisher class itself is not installed on the assembly host, so the wrapper class amsart is used and the equation counter is still tied to the section counter as in the source.\par\medskip
\noindent\textit{Editorial note.} The retained bibliography carries exactly the one work cited by the retained text (Nazarov, Quantum Berezinian and the classical Capelli identity), copied verbatim from the source's own reference list.\par\medskip
\noindent\textit{Editorial scope.} One original result: Proposition 3.1 (label zu) of the source with its complete author proof at source lines 888-977, together with the author's own definition of the super Yangian that the proof uses. The already-counted Theorem 2.8 of the same source, the free-generation theorem of the same section and the quantum Berezinian shift theorem are not counted here and are not used as premises.\par\medskip
\noindent\textit{Allowed standard prerequisites (editorial).}\begin{itemize}
\item Super Yangian and its RTT presentation. For the general linear Lie superalgebra gl\_\{M|N\} over C the Yangian Y(gl\_\{M|N\}) is the associative Z\_2-graded algebra with generators T\_ij\textasciicircum{}(r) for r at least 1, whose generating series T\_ij(u) satisfy the RTT relation; the inverse matrix series is written T'\_ij(u) and the antipode is written S. The author sets all of this up at source lines 169-320, which are retained in the material, so nothing in the counted proof depends on an outside convention.. Fixes the algebra, its generators, the series T\_ij(u) and T'\_ij(u) and the antiautomorphism tau from which the operator Q is built.
\item Formal power series in u\textasciicircum{}\{-1\} and the rational R-matrix. All series are formal in the variable u\textasciicircum{}\{-1\} with coefficients in the Yangian, and the rational R-matrix R(u) together with its twisted inverse R-hat\_S(u) = 1 + Q(u-M+N)\textasciicircum{}\{-1\} is used through the identities recorded at source lines 249-308, which are retained in the material.. Supplies the identities (3.0), (3.3), (3.52) and (pop) that the counted proof quotes when it forms Q, R-hat\_S and the three-leg computation.
\end{itemize}\medskip
\section{Definition of the Yangian}\label{sec:1}

Throughout this article we will use the following general conventions.
Let $\mathrm{A}$ and $\mathrm{B}$ be any two associative
$\ZZ_{2}$-graded algebras. Their tensor product
$\mathrm{A}\ot\mathrm{B}$ is also
an associative $\ZZ_2$-graded algebra such that for any homogeneous
elements
$X,X^\prime\in\mathrm{A}$ and $Y,Y^\prime\in\mathrm{B}$
\begin{gather}\label{conv1}
(X\ot Y) (X^\prime\ot Y^\prime)=X X^\prime\ot Y Y^\prime (-1)^{\deg X^\prime\deg Y},\\
\label{conv2}
\deg (X\ot Y)=\deg X+\deg Y .
\end{gather}
For any $\ZZ_2$-graded modules
$U$ and $V$ over $\mathrm{A}$ and $\mathrm{B}$ respectively,
the vector space $U\ot V$ is a~$\ZZ_2$-graded module over
$\mathrm{A}\ot\mathrm{B}$ such that for any homogeneous elements
$x\in U$ and $y\in V$
\begin{gather}\label{XY}
(X\ot Y) ( x\ot y ) =X x\ot Y y (-1)^{\deg x \deg Y},\\
\label{xy}
\deg ( x\ot y )=\deg x+\deg y .
\end{gather}
A homomorphism $\al\colon \mathrm{A}\to\mathrm{B}$ is a linear map
such that $\al (X X^\prime)=\al (X) \al (X^\prime)$ for all
$X,X^\prime\in\mathrm{A} $. But an antihomomorphism
$\be\colon \mathrm{A}\to\mathrm{B}$ is a linear map
such that for all homogeneous $X,X^\prime\in\mathrm{A}$
\begin{gather}\label{bb}
\be (X X^\prime )=\be (X^\prime) \be (X) (-1)^{\deg X\deg X^\prime}.
\end{gather}

If $\mathrm{A}$ is unital,
let $\iota_{h}$ be
its embedding into the tensor product $\mathrm{A}^{\ns\ot n}$
as the $h $-th tensor factor:
\begin{gather*}
\iota_{h}(X)=1^{\ot (h-1)}\ot X\ot1^{\ot (n-h)}\qquad\text{for}\quad h=1,\dots, n .
\end{gather*}
Here $n$ can be any positive integer.
We will also use various embeddings of the algebra
$\mathrm{A}^{\ns\ot m}$ into~$\mathrm{A}^{\ns\ot n}$
for any $m=1,\dots, n$.
For any choice of pairwise distinct indices $h_1,\dots, h_m\in\{1,\dots, n \}$
and of an element $X\in\mathrm{A}^{\ns\ot m}$ of the form
$X=X^{(1)}\ot\cdots\ot X^{(m)}$ we will denote
\begin{gather*}
X_{h_1\dots h_m}=\iota_{h_1}\big(X^{(1)}\big) \cdots \iota_{h_m}\big(X^{(m)}\big) \in\mathrm{A}^{\ns\ot n}.
\end{gather*}
We will then extend the notation $X_{h_1\dots h_m}$ to all elements
$X\in\mathrm{A}^{\ns\ot m}$ by linearity.

Now let the indices $i$, $j$ run through $1,\dots, M+N $.
We will always write $\bi=0$ if $1\le i\le M$ and
$\bi=1$ if $M<i\le M+N $. Consider the
$\ZZ_{2}$-graded vector space $\CMN $.
Let $e_i\in\CMN$ be an element of the standard basis.
The $\ZZ_{2}$-grading on $\CMN$ is defined so that $\deg e_i=\bi $.
Let $E_{ij}\in\End\CMN$ be the standard matrix unit, so that
$E_{ij} e_k=\de_{jk} e_i $. The
associative algebra $\End\CMN$
is $\ZZ_{2}$-graded so that $\deg E_{ij}=\bi+ \bj $.

For any $n$ we can identify
the tensor product $\big(\!\End\CMN\big)^{\ot n}$ with the algebra
$\End\!\big(\big(\CMN\big)^{\!\ot n}\big)$ acting on the vector space $\big(\CMN\big)^{\!\ot n}$
by repeatedly using the conventions~\eqref{XY} and~\eqref{xy}.

Let us introduce the {\it Yangian} of the Lie superalgebra
$\glMN$. This is the complex associative
unital $\ZZ_{2}$-graded algebra $\YMN$ with the countable set of generators
\begin{gather}\label{Tijr}
T^{(r)}_{ij},
\qquad\text{where}\quad r=1,2, \dots \quad\text{and}\quad i ,j=1,\dots, M+N .
\end{gather}
The $\ZZ_{2}$-grading on the algebra $\YMN$
is determined by setting $\deg T_{ij}^{(r)}=\bi+\bj$ for $r\ge1$.
To write down defining relations for these
generators we will employ the series
\begin{gather}\label{3.0}
T_{ij}(u)=\de_{ij}\cdot1+T_{ij}^{(1)} u^{-1}+T_{ij}^{(2)} u^{-2}+\cdots
\end{gather}
in a formal variable $u$ with coefficients from $\YMN $.
Then for all possible indices $i$, $j$, $k$, $l${\samepage
\begin{gather}\label{3.1}
(u-v)[ T_{ij}(u) ,T_{kl}(v) ](-1)^{\bi \bk + \bi \bl + \bk \bl}=T_{kj}(u) T_{il}(v)-T_{kj}(v) T_{il}(u),
\end{gather}
where $v$ is another formal variable.
The square brackets here stand for the supercommutator.
Notice that the series denoted by $T_{ij}(u)$
here and in \cite{N1} differ by the scalar factor
$(-1)^{( \bi+1 ) \bj}$.}

We will also use the following matrix form of the defining relations~\eqref{3.1}.
Take the element
\begin{gather}\label{pop}
P = \sum_{i,j=1}^{M+N} E_{ij}\ot E_{ji} (-1)^{\bj}
\in\big(\End\CMN\big)^{\ot 2} .
\end{gather}
This element acts on the vector space $\big(\CMN\big)^{\ot 2}$ so that
$e_i\ot e_{j}\mapsto e_j\ot e_i {(-1)}^{\bi \bj}$.
Here we identify the algebra $\big(\End\CMN\big)^{\ot 2}$ with the algebra
$\End\big(\big(\CMN\big)^{\ot 2}\big)$ by using~\eqref{XY}.

For any $n$ let $\Sym_n$ be the symmetric group acting
on the set $\{1,\dots, n\}$ by permutations. For each $m=1,\dots, n-1$
denote by $\si_m$ the element of $\Sym_n$ exchanging~$m$ and $m+1 $.
The group $\Sym_n$ also acts on the vector space $\big(\CMN\big)^{\ot n}$.
This action is defined by the assignment
$\si_m\mapsto P_{m,m+1}$ for each~$m $.
Here we identify the algebra
$\big(\End\CMN\big)^{\ot n}$ with $\End\big(\big(\CMN\big)^{\ot n}\big)$ via~\eqref{XY},~\eqref{xy}.

The rational function $R(u)=1-P u\1$ with values in the algebra
$\big(\End\CMN\big)^{\ot 2}$ is called the {\it Yang $R$-matrix}.
It satisfies the {\it Yang--Baxter equation}
in the algebra $\big(\End\CMN\big)^{\ot 3}(u,v,w)$
\begin{gather}\label{3.6}
R_{12}(u-v) R_{13}(u-w) R_{23}(v-w)=R_{23}(v-w) R_{13}(u-w) R_{12}(u-v) .
\end{gather}
Since $P^2=1 $, we also have the relation
\begin{gather}\label{3.52}
R(-u) R(u)=1-u^{-2} .
\end{gather}

Now combine all the series \eqref{3.0} into the single element
\begin{gather}\label{tu}
T(u)=\sum_{i,j=1}^{M+N} E_{ij}\ot T_{ij}(u)
\in\big(\End\CMN\big)\ot\YMN \big[\big[u^{-1}\big]\big] .
\end{gather}
For any $n$ and any $p=1,\dots, n$ we will denote
\begin{gather}
T_p(u)=\iota_p\ot\id ( T(u)) \in\big(\End\CMN\big)^{\ot n}\ot\YMN \big[\big[u\1\big]\big] .\label{3.22}
\end{gather}
By using this notation for $n=2$
the relations \eqref{3.1} can be rewritten as
\begin{gather}\label{3.3}
(R (u-v)\ot1) T_1(u) T_2(v)=T_2(v) T_1(u) (R (u-v)\ot1) .
\end{gather}
Namely, after multiplying each side of \eqref{3.3} by $u-v$ it becomes
a relation of series in $u$, $v$ with coefficients in %the algebra
$\big(\End\CMN\big)^{\ot 2}\ot\YMN$
equivalent to the collection of all the relations~\eqref{3.1}.

\begin{Proposition}\label{m}
An antiautomorphism of $\YMN$ can be defined by the assignment
\begin{gather}\label{M}
T_{ij}(u)\mapsto T_{ij}(-u) .
\end{gather}
\end{Proposition}

\begin{proof}Due to the convention \eqref{conv1}, by using
the notation \eqref{3.22} for $n=2$ we get
\begin{gather}\label{TT}
T_1(u) T_2(v)=\sum_{i,j,k,l=1}^{M+N}
E_{ij}\ot E_{kl}\ot T_{ij}(u) T_{kl}(v)
{(-1)}^{( \bi+ \bj )( \bk+ \bl )} ,\\
\label{MM}
T_2(-v) T_1(-u)=\sum_{i,j,k,l=1}^{M+N}
E_{ij}\ot E_{kl}\ot T_{kl}(-v) T_{ij}(-u) .
\end{gather}
By using the convention \eqref{bb}, the antihomomorphism property of \eqref{M} follows from the relation
\begin{gather*}
(R (u-v)\ot1) T_2(-v) T_1(-u)=
T_1(-u) T_2(-v) (R (u-v)\ot1),
\end{gather*}
which is obtained from~\eqref{3.3} by using~\eqref{3.52}.
The antihomomorphism~\eqref{M} is clearly involutive and therefore bijective.
\end{proof}

For all indices $i$, $j$ define the series $\TP_{ij}(u)$
by using the element inverse to \eqref{tu} so that
\begin{gather}\label{tunat}
T(u)^{-1}=\sum_{i,j=1}^{M+N} E_{ij}\ot\TP_{ij}(u) .
\end{gather}

\begin{Proposition}\label{s}
An antiautomorphism of $\YMN$ can be defined by the assignment
\begin{gather}\label{S}
T_{ij}(u)\mapsto\TP_{ij}(u) .
\end{gather}
\end{Proposition}

Further, let $\tau$ be the antiautomorphism of $\End\CMN$
defined by the assignment
\begin{gather*}
E_{ij}\mapsto E_{ji} {(-1)}^{\bi ( \bj + 1)} .
\end{gather*}
Then by the definition \eqref{tu} we have
\begin{gather*}
\tau \ot\id ( T(u) )=
\sum_{i,j=1}^{M+N} E_{ji}\ot T_{ij}(u)
{(-1)}^{\bi ( \bj + 1)}=
\sum_{i,j=1}^{M+N} E_{ij}\ot T_{ji}(u)
{(-1)}^{\bj ( \bi + 1)} .
\end{gather*}

\begin{Proposition}
An antiautomorphism of %the algebra
$\YMN$ can be defined by the assignment
\begin{gather}
\label{T}
T_{ij}(u)\mapsto T_{ji}(u)
{(-1)}^{\bj ( \bi + 1)} .
\end{gather}
\end{Proposition}

\begin{proof}
Observe that $(\tau\ot\tau)(P)=P$ and hence
\begin{gather}
\label{Rtt}
(\tau\ot\tau)(R(u-v))=R(u-v) .
\end{gather}
Therefore the antihomomorphism property of \eqref{T}
follows from the relation which is obtained by applying
$\tau\ot\tau\ot\id$ to both sides of\eqref{3.3},
see the proofs of Propositions \ref{m} and \ref{s} above.
To prove the bijectivity of
the antihomomorphism \eqref{T}, observe
that its square is given by %the assignment
\begin{gather*}
T_{ij}(u)\mapsto T_{ij}(u)
{(-1)}^{\bi+ \bj} .
\end{gather*}
In particular, the square is an automorphism of the $\ZZ_2$-graded
algebra $\YMN $.
\end{proof}

Put
$\TS(u)=\tau \ot\id \big( T(u)^{-1} \big) $.
Then by \eqref{tunat}
\begin{gather}\label{tsu}
\TS(u)=
\sum_{i,j=1}^{M+N} E_{ji}\ot\TP_{ij}(u)
{(-1)}^{\bi ( \bj + 1)}=
\sum_{i,j=1}^{M+N} E_{ij}\ot\TP_{ji}(u)
{(-1)}^{\bj ( \bi + 1)} .
\end{gather}

\section{Centre of the Yangian}\label{sec:2}

Here we will give a description of the centre of the
algebra $\YMN $. An element of $\YMN$ is called
\emph {central} if it supercommutes with each element of $\YMN $.
However, we will see that the central elements of $\YMN$ have
the $\ZZ_{2}$-degree $0 $.
Hence they commute with each
element of~$\YMN$ in the usual sense.
Our description comes from computing the square of the antipodal mapping $\S$ of $\YMN $.
Another description of the centre of~$\YMN $ will be
given in the next section. In that section we will also establish a
correspondence between the two descriptions.


\begin{Proposition}\label{zu} There is a formal power series $Z(u)$ in $u^{-1}$ with coefficients in
the centre of~$\YMN$ and with
leading term~$1$ such that for all indices $i$ and $j$
\begin{gather}\label{zetav}
\sum_{k=1}^{M+N}T_{kj}(u+M-N) \TP_{ik}(u) =\delta_{ij} Z(u) ,\\
\label{zetau}
\sum_{k=1}^{M+N}
\TP_{kj}(u) T_{ik}(u+M-N) =\delta_{ij} Z(u) .
\end{gather}
\end{Proposition}

\begin{proof}By using the definition \eqref{pop} introduce the element of
the algebra $\big(\End\CMN\big)^{\ot 2}$
\begin{gather}\label{Q}
Q=( \id\ot\tau ) ( P )=\sum_{i,j=1}^{M+N}E_{ij}\ot E_{ij} {(-1)}^{\bi \bj} .
\end{gather}
The image of the action of $Q$ on $\big(\CMN\big)^{\ot 2}$ is one dimensional and is spanned by the vector
\begin{gather*}
\sum_{i=1}^{M+N}e_{i}\ot e_{i} .
\end{gather*}
Here we regard $Q$ as an element of the algebra
$\End\big(\big(\CMN\big)^{\ot 2}\big)$
by identifying this algebra with $\big(\End\CMN\big)^{\ot 2}$ via~\eqref{XY}.
We also have
$Q^{2}=(M-N) Q $. By using the latter relation we get
\begin{gather}\label{rutinv}
(( \id\ot\tau ) ( R(u) ))^{-1}=\big( 1-Q u\1\big)^{-1}=
1+Q ( u-M+N )^{-1} .
\end{gather}
The rational function of the variable~$u$
given by the equalities \eqref{rutinv} will be denoted by $\RS(u) $.

Now multiply both sides of the relation~\eqref{3.3}
by $T_2^{-1}(v)$ on the left and right, and then apply~$\tau$ relative
to the second tensor factor $\End\CMN$ in $\big(\End\CMN\big)^{\ot2}\ot\YMN $.
Multiplying the resulting relation on the left and right by
$\RS(u-v)\ot 1$ yields
\begin{gather}\label{trater}
\big(\RS(u-v)\ot1\big) T_1(u) \TS_2(v)=\TS_2(v) T_1(u)(\RS(u-v)\ot1) .
\end{gather}
Multiplying the latter relation by $u-v-M+N$ and then setting
$u=v+M-N$ we get
\begin{gather}\label{qresi}
(Q\ot1) T_1(v+M-N) \TS_2(v)= \TS_2(v) T_1(v+M-N) (Q\ot1) ,
\end{gather}
see \eqref{rutinv}.
As the image of $Q$ in $\big(\CMN\big)^{\ot 2}$
is one dimensional, either side of the relation~\eqref{qresi}
must be equal to $Q\ot Z(v)$ where~$Z(v)$
is a certain power series in~$v^{-1}$
with coefficients from~$\YMN $.
The equality
of the left-hand side and of the right-hand side of~\eqref{qresi}
to $Q\ot Z(v)$
is equivalent respectively to~\eqref{zetav} and~\eqref{zetau}.
We just need to replace the variable~$v$ in~\eqref{qresi} by~$u $.

It is immediate from \eqref{3.0} and \eqref{zetav} that
every coefficient of the series $Z(u)$ has
$\ZZ_{2}$-degree $0$ and that its leading
term is~$1 $.
Let us prove that all these coefficients
are central in $\YMN $.

To this end we will work with the elements \eqref{3.22}
where $n=3 $. By using \eqref{3.3} and \eqref{trater},
\begin{gather}
\big(\RS_{13}(u-v) R_{12}(u-v-M+N)\ot1\big)
T_1(u) T_2(v+M-N) \TS_3(v)
\notag
\\
\qquad{}=\big(\RS_{13}(u-v)\ot1\big)
T_2(v+M-N) T_1(u) \TS_3(v)
(R_{12}(u-v-M+N)\ot1)
\notag
\\
\qquad{} =T_2(v+M-N) \TS_3(v) T_1(u)
\big(\RS_{13}(u-v) R_{12}(u-v-M+N)\ot1\big) .\label{3.666666}
\end{gather}
On the other hand, due to \eqref{3.52} we have the identity in the algebra
$\big(\End\CMN\big)^{\ot 3}( u )$
\begin{gather*}
R_{13}(-u) P_{23} R_{12}(u)=\big( 1-u^{-2} \big) P_{23} .
\end{gather*}
By applying to it the antiautomorphism~$\tau$
relative to the third tensor factor of $\big(\End\CMN\big)^{\ot 3}$
and then using the definition~\eqref{rutinv} we get
\begin{gather}\label{QR}
Q_{23} \RS_{13}(u+M-N) R_{12}(u)=\big( 1-u^{-2} \big) Q_{23} .
\end{gather}
Therefore multiplying the first and third lines of the display~\eqref{3.666666} by $Q_{23}\ot1$ on the left yields
\begin{gather*}
\big( 1-(u-v-M+N)^{-2} \big)
T_1(u) (Q_{23}\ot Z(v))\\
\qquad{} = (Q_{23}\ot Z(v)) T_1(u)\big( 1-(u-v-M+N)^{-2} \big) .
\end{gather*}
We just need to replace the variable $u$ in~\eqref{QR} by~$u-v-M+N$
and use the relation~\eqref{qresi}.
The last relation implies that any generator~$T_{ij}^{(r)}$
commutes with every coefficient of~$Z(v)$.
\end{proof}

\begin{thebibliography}{99}
\bibitem{N1}
Nazarov M.L., Quantum {B}erezinian and the classical {C}apelli identity,
 \href{https://doi.org/10.1007/BF00401646}{\textit{Lett. Math. Phys.}} \textbf{21} (1991), 123--131.
\end{thebibliography}
\end{document}
